\begin{frame}[plain]
\wrsTitlePage{LearniBridge}{What correction makes a cached computation behave like the skipped full computation?}{Paper explainer \textperiodcentered{} arXiv:2606.26778 \textperiodcentered{} ICML 2026 \textperiodcentered{} an independent redrawing of the mechanism, not a reproduction of the paper's figures}{Huang, Chen, Shen \& Zhang \textperiodcentered{} ICML 2026}
\note{This is a method explainer, not a leaderboard review. The whole deck answers one question: what does the correction have to look like for cache reuse to survive long skips?}
\end{frame}

\begin{frame}{Feature caching buys speed by skipping blocks}
\begin{columns}[T]
\column{0.56\textwidth}
\wrsLead{A Diffusion Transformer evaluates a deep network at every denoising step. Feature caching reuses a representation computed at an earlier step so later steps can skip expensive blocks.}\par\vspace{0.6em}
\wrsSection{Why it is hard}
\begin{itemize}
  \item The reused feature is stale. It was computed for a different timestep, so it carries the wrong content.
  \item The same approximation is applied at every skipped step, and the error compounds along the trajectory.
  \item The problem grows exactly where the win is. The longer the cache interval, the larger the drift and the larger the speedup.
\end{itemize}
\column{0.4\textwidth}
\wrsSection{Constraints}
\begin{itemize}
  \item The backbone stays as it is. No architectural change and no retraining of the base model.
  \item The calibration budget must be small, a handful of prompts rather than a finetune.
  \item Skipped steps must stay cheap, otherwise the speedup is spent on the fix.
\end{itemize}
\end{columns}
\note{Establish the trade-off before any architecture. Caching is attractive because it is model-agnostic and needs no training, and that is exactly the property a fix should preserve.}
\end{frame}

\begin{frame}{Existing methods forecast from history}
\begin{columns}[T]
\column{0.31\textwidth}
\wrsSection{A. Intuition}
\wrsSmall{Nearby steps produce similar features, so an earlier feature stands in well.}
\column{0.31\textwidth}
\wrsSection{B. Mechanism}
\wrsSmall{Cache the block output at a reference step, then approximate skipped steps from history: reuse (DeepCache, FORA) or Taylor extrapolation (TaylorSeer).}
\column{0.31\textwidth}
\wrsSection{C. Where it fails}
\begin{itemize}
  \item A fixed rule cannot track how the representation changes away from the reference step.
  \item TaylorSeer assumes smooth, high-order differentiable trajectories.
  \item At high acceleration the error accumulates.
\end{itemize}\vspace{0.3em}
\begin{wrsband}{interpretation}{Relation to us}Keep the schedule; replace the rule with a learned, structured correction.\end{wrsband}
\end{columns}
\note{Keep the baselines compressed to one idea each. The point is not the catalogue but the shared assumption, a fixed rule over historical features.}
\end{frame}

\begin{frame}{Change the question, learn the correction}
\begin{columns}[T]
\column{0.3\textwidth}
\wrsSection{The usual question}
\wrsSmall{Predict the future feature from historical features with a fixed rule, and hope the prediction is close.}
\column{0.3\textwidth}
\wrsSection{LearniBridge asks}
\wrsSmall{What correction is required for the cached computation to reproduce the skipped full computation?}
\column{0.32\textwidth}
\wrsSection{The idea}
\wrsBody{Model the required change directly and constrain it. The cached input times a low-rank correction approximates the future full-computation output.}
\end{columns}
\vspace{0.6em}\wrsBody{If the required correction has exploitable structure, a few calibration prompts are enough and the correction is nearly free at inference.}
\note{This is the pivot of the paper. Forecasting works in feature space and guesses forward; LearniBridge asks what change in the computation is needed to land on the true output.}
\end{frame}

\begin{frame}{What correction is actually required}
\begin{center}
\begin{tikzpicture}[x=1cm,y=1cm]
  \node[wrsnode] (lb-cached) at (1.38, 2.25) {$X$};
  \node[wrslabel, below=0.02cm of lb-cached] {cached input};
  \node[wrsnode] (lb-reuse) at (4.49, 4.75) {$\tilde{F}$};
  \node[wrslabel, below=0.02cm of lb-reuse] {output if we just reuse it};
  \node[wrsnode] (lb-target) at (7.93, 2.50) {$F$};
  \node[wrslabel, below=0.02cm of lb-target] {true output};
  \draw[wrsreference] (lb-cached) -- node[midway, above, wrslabel] {reuse} (lb-reuse);
  \draw[wrsvector] (lb-reuse) -- node[midway, above, wrslabel] {residual E} (lb-target);
\end{tikzpicture}
\vspace{0.2em}\wrsLegendVector{residual E: what the correction must explain}\hspace{1.6em}\wrsLegendReference{reuse path, with no correction}
\end{center}
\begin{center}\wrsQuiet{Over many prompts the paper writes the residuals as a matrix E and the cached inputs as a matrix X, and models E $\approx$ $\Delta$W \textperiodcentered{} X. The geometry above is in feature space; the correction $\Delta$W itself acts on layer weights.}\end{center}
\note{Keep the three objects on screen. X is one cached input, F\textasciitilde{} is what reuse returns, F is what full computation would have returned. The correction is defined by that gap, and the next slide asks what shape the gap has.}
\end{frame}

\begin{frame}{Diagnostic D1: is the required correction low-rank?}
\wrsDiagId{D1}\quad\wrsBody{Does the optimal calibration update have exploitable structure?}\hfill\wrsTests{C1}
\vspace{0.6em}
\begin{columns}[T]
\column{0.56\textwidth}
\begin{wrsband}{measurement}{Measurement}\wrsMono{Singular values of X (100 prompts) decay rapidly. Fit E $\approx$ $\Delta$W \textperiodcentered{} X: $\Delta$W = E \textperiodcentered{} pinv(X), rank($\Delta$W) $\le$ rank(X) $\le$ r}\end{wrsband}
\begin{wrsband}{observation}{Observation}Spectral energy concentrates in a few principal directions, so the least-squares correction is low rank.\end{wrsband}
\begin{wrsband}{interpretation}{Interpretation}A low-rank parameterization can represent the required correction; no full-rank update is needed.\end{wrsband}
\column{0.4\textwidth}
\begin{wrsband}{can}{Can conclude}Under the paper's model E $\approx$ $\Delta$W \textperiodcentered{} X, the optimal $\Delta$W is low-rank in the authors' setting.\end{wrsband}
\begin{wrsband}{cannot}{Cannot conclude}That diffusion features are low-rank in general, or that the bound holds for every layer and model.\end{wrsband}
\begin{wrsband}{alternative}{Alternative targeted}The decay may partly reflect the chosen layer and prompt pool; this is empirical, not universal.\end{wrsband}
\end{columns}
\vfill\wrsCitation{Huang et al., ICML 2026, §3.2 and Figure 1}
\note{Walk the bands in order. The measurement is the SVD and the least-squares solution; the observation is the concentration of spectral energy; the interpretation is that low rank is a usable parameterization. The bound rank($\Delta$W) $\le$ rank(X) is a consequence of the paper's model, not a theorem about diffusion models.}
\end{frame}

\begin{frame}{Diagnostic D2: is the correction shared across prompts?}
\wrsDiagId{D2}\quad\wrsBody{Does the required correction depend on the prompt?}\hfill\wrsTests{C1}
\vspace{0.6em}
\begin{columns}[T]
\column{0.56\textwidth}
\begin{wrsband}{measurement}{Measurement}\wrsMono{Fit $\Delta$W on each of 20 disjoint prompt groups; compare principal subspaces pairwise: the angles are small.}\end{wrsband}
\begin{wrsband}{observation}{Observation}Different prompt groups need nearly the same correction subspace.\end{wrsband}
\begin{wrsband}{interpretation}{Interpretation}The pattern is stable enough that calibration from a few prompts transfers.\end{wrsband}
\column{0.4\textwidth}
\begin{wrsband}{can}{Can conclude}In the tested setting, subspaces from disjoint prompt groups coincide.\end{wrsband}
\begin{wrsband}{cannot}{Cannot conclude}That every prompt needs the same correction, or that invariance survives a domain shift.\end{wrsband}
\begin{wrsband}{alternative}{Alternative targeted}The groups came from one prompt pool, which may inflate their similarity.\end{wrsband}
\end{columns}
\vfill\wrsCitation{Huang et al., ICML 2026, Figure 2}
\note{This is the evidence behind the 3-5 prompt calibration budget. Be precise: small principal-subspace angles, not identical corrections, and measured on a shared prompt pool.}
\end{frame}

\begin{frame}{Low rank suggests a LoRA bridge on the final block}
\begin{center}
\resizebox{\ifdim\width>\linewidth\linewidth\else\width\fi}{!}{%
\begin{tikzpicture}[node distance=2.80cm]
  \node[wrstoken] (stage-0) {cached input};
  \node[wrslearned, right=of stage-0] (stage-1) {final block + $\Delta$W = BA};
  \node[wrslatent, right=of stage-1] (stage-2) {calibrated output};
  \draw[wrsarrow] (stage-0) -- (stage-1);
  \draw[wrsarrow] (stage-1) -- (stage-2);
  \node[wrsdetail, below=0.45cm of stage-0] {the final-block input recorded at the last full step};
  \node[wrsdetail, below=0.45cm of stage-1] {LoRA on Q, K, V, O, W1, W2 of block L; the base weights stay frozen};
  \node[wrsdetail, below=0.45cm of stage-2] {approximates the full computation at the skipped step};
\end{tikzpicture}}
\end{center}
\vspace{0.3em}\begin{center}\wrsQuiet{Low rank is not an extra constraint here: LoRA is exactly a low-rank parameterization, $\Delta$W = BA with r far below the layer dimensions. The only trainable parameters are the two small factors A and B.}\end{center}
\note{Make the design follow from the evidence: the analysis says the update is low rank, LoRA is the natural way to parameterize a low-rank update, so the adapter is a consequence rather than a fashion choice.}
\end{frame}

\begin{frame}{Calibration: one full pass, then only the adapters train}
\begin{columns}[T]
\column{0.5\textwidth}
\wrsSection{A. Pre-calibration pass (full computation)}
\begin{itemize}
  \item \textbf{Record input.} At each full-computation step t, keep the final-block input.
  \item \textbf{Record target.} At the steps that will be skipped, keep the true final-block output.
\end{itemize}
\column{0.46\textwidth}
\wrsSection{B. LoRA finetune (only the adapters train)}
\wrsBody{The LoRA-augmented final block must turn the cached input into that recorded target, while every base weight stays frozen.}\par\vspace{0.4em}
\wrsTiny{Budget reported in the paper: 3-5 calibration prompts, small rank, no backbone finetuning.}
\end{columns}
\note{Be honest about the cost: calibration needs a pre-calibration pass at full computation, so the training data is not free, it is just small. The ablation shows performance saturating beyond about five prompts.}
\end{frame}

\begin{frame}{Inference: skip the blocks, run the bridge}
\begin{center}
\resizebox{\ifdim\width>\linewidth\linewidth\else\width\fi}{!}{%
\begin{tikzpicture}[node distance=2.80cm]
  \node[wrscache] (stage-0) {full step};
  \node[wrsstage, right=of stage-0] (stage-1) {skipped step};
  \node[wrslearned, right=of stage-1] (stage-2) {LoRA final block};
  \draw[wrsarrow] (stage-0) -- (stage-1);
  \draw[wrsarrow] (stage-1) -- (stage-2);
  \node[wrsdetail, below=0.45cm of stage-0] {every N steps: run all blocks and cache the final-block input};
  \node[wrsdetail, below=0.45cm of stage-1] {blocks g1 to gL-1 are bypassed; no computation is spent there};
  \node[wrsdetail, below=0.45cm of stage-2] {only the last block runs, corrected by the low-rank bridge};
\end{tikzpicture}}
\end{center}
\vspace{0.3em}\begin{center}\wrsQuiet{This is the whole acceleration: full computation every N steps, one calibrated block in between. The bridge is not run through every block.}\end{center}
\note{Say the negative explicitly, because it is a common misreading: the adapters live in the final block only, and the earlier blocks do not run at all at skipped steps.}
\end{frame}

\begin{frame}{Text-to-image: FLUX.1-dev at about 6x}
\begin{center}\scriptsize
\setlength{\tabcolsep}{3pt}
\renewcommand{\arraystretch}{1.15}
\begin{tabular}{lccccc}
\toprule
Method & \textbf{Speed} & \textbf{ImageReward} & \textbf{CLIP} & \textbf{PSNR} & \textbf{SSIM} \\
\midrule
ToCA (N=12) & 5.87 & 0.70 & 0.78 & 26.48 & 0.59 \\
TeaCache (delta=1.4) & 5.99 & 0.73 & 0.80 & 26.58 & 0.63 \\
TaylorSeer (N=8, O=2) & 6.10 & 0.82 & 0.80 & 26.82 & 0.68 \\
\textbf{\textcolor{wrsBlue}{LearniBridge (N=8)}}\newline{\tiny\color{wrsMuted}(ours)} & \textbf{\textcolor{wrsBlue}{6.17}} & \textbf{\textcolor{wrsBlue}{0.83}} & \textbf{\textcolor{wrsBlue}{0.82}} & \textbf{\textcolor{wrsBlue}{28.35}} & \textbf{\textcolor{wrsBlue}{0.69}} \\
\bottomrule
\end{tabular}
\end{center}
\vspace{0.25em}\wrsTiny{higher is better}
\vspace{0.25em}\wrsQuiet{200 DrawBench prompts. At a slightly higher speedup, LearniBridge keeps the best ImageReward, CLIP, PSNR and SSIM of this group. PSNR and SSIM measure fidelity to the unaccelerated output, not perceptual quality on their own.}
\vfill\wrsCitation{Huang et al., ICML 2026, Table 1}
\note{Read the table as a trade curve, not as a win list: the baselines are allowed different settings. The robust signal is that the ranking holds as the ratio grows, which is what the paper argues.}
\end{frame}

\begin{frame}{Text-to-video: WAN 2.1-1.3B on VBench}
\begin{center}\scriptsize
\setlength{\tabcolsep}{3pt}
\renewcommand{\arraystretch}{1.15}
\begin{tabular}{lcccc}
\toprule
Method & \textbf{Speed} & \textbf{VBench} & \textbf{PSNR} & \textbf{SSIM} \\
\midrule
EasyCache (delta=0.13) & 3.27 & 79.61 & 13.77 & 0.47 \\
TaylorSeer (N=4, O=2) & 3.46 & 79.17 & 14.84 & 0.45 \\
TeaCache (delta=0.3) & 3.98 & 79.23 & 12.82 & 0.47 \\
\textbf{\textcolor{wrsBlue}{LearniBridge (N=5)}}\newline{\tiny\color{wrsMuted}(ours)} & \textbf{\textcolor{wrsBlue}{4.10}} & \textbf{\textcolor{wrsBlue}{80.51}} & \textbf{\textcolor{wrsBlue}{16.32}} & \textbf{\textcolor{wrsBlue}{0.63}} \\
\textbf{\textcolor{wrsBlue}{LearniBridge (N=4)}}\newline{\tiny\color{wrsMuted}(ours)} & \textbf{\textcolor{wrsBlue}{3.49}} & \textbf{\textcolor{wrsBlue}{81.21}} & \textbf{\textcolor{wrsBlue}{21.26}} & \textbf{\textcolor{wrsBlue}{0.82}} \\
\bottomrule
\end{tabular}
\end{center}
\vspace{0.25em}\wrsTiny{higher is better}
\vspace{0.25em}\wrsQuiet{946 prompts, five samples each. At 4.10x the bridge keeps VBench 1.28 points above the best baseline at a comparable ratio, and at 3.49x it stays close to the unaccelerated model. HunyuanVideo shows the same pattern up to 5.75x in the paper's Table 2.}
\vfill\wrsCitation{Huang et al., ICML 2026, Table 3}
\note{Two LearniBridge rows on purpose: one at a higher ratio than every baseline and one at a moderate ratio. Point out that PSNR and SSIM here measure agreement with the unaccelerated video of the same model.}
\end{frame}

\begin{frame}{Our runs: N=5 and N=7 (CLIP, IR, HPSv2.1, PSNR, SSIM, speed)}
\begin{center}\scriptsize
\setlength{\tabcolsep}{3pt}
\renewcommand{\arraystretch}{1.15}
\begin{tabular}{lcccccc}
\toprule
Method & \textbf{CLIP} & \textbf{IR} & \textbf{HPSv2.1} & \textbf{PSNR} & \textbf{SSIM} & \textbf{Speed} \\
\midrule
Original FLUX & 32.65 & 1.00 & 30.40 & $\infty$ & 1 & 1.00x \\
Cache thường (N=5) & 32.63 & 0.94 & 29.44 & 14.70 & 0.55 & 4.63x \\
Global-Local (N=5) & 32.79 & 1.01 & 30.00 & 15.21 & 0.57 & 4.60x \\
\textbf{\textcolor{wrsBlue}{LearniBridge-100 (N=5)}}\newline{\tiny\color{wrsMuted}(paper method)} & \textbf{\textcolor{wrsBlue}{32.66}} & \textbf{\textcolor{wrsBlue}{0.98}} & \textbf{\textcolor{wrsBlue}{29.53}} & \textbf{\textcolor{wrsBlue}{15.07}} & \textbf{\textcolor{wrsBlue}{0.55}} & \textbf{\textcolor{wrsBlue}{4.49x}} \\
Cache thường (N=7) & 32.28 & 0.83 & 28.24 & 14.14 & 0.52 & 5.66x \\
Global-Local (N=7) & 32.63 & 0.92 & 29.10 & 14.73 & 0.55 & 5.60x \\
\textbf{\textcolor{wrsBlue}{LearniBridge-100 (N=7)}}\newline{\tiny\color{wrsMuted}(paper method)} & \textbf{\textcolor{wrsBlue}{32.76}} & \textbf{\textcolor{wrsBlue}{0.92}} & \textbf{\textcolor{wrsBlue}{28.44}} & \textbf{\textcolor{wrsBlue}{14.49}} & \textbf{\textcolor{wrsBlue}{0.54}} & \textbf{\textcolor{wrsBlue}{5.45x}} \\
\bottomrule
\end{tabular}
\end{center}
\vspace{0.25em}\wrsTiny{higher is better}
\vspace{0.25em}\wrsQuiet{Fidelity metrics PSNR/SSIM compare against the unaccelerated Original FLUX output, so they fall for every accelerator including the best one; CLIP/IR/HPSv2.1 are the quality signals. LPIPS is omitted here for width (it tracks PSNR/SSIM).}
\note{Baseline naming is shown as given in our run logs. Read the rows by N: at the same cache interval, Global-Local is strongest on fidelity while LearniBridge-100 stays closest on CLIP at the longer interval.}
\end{frame}

\begin{frame}{Our runs: N=10 and N=25 (quality collapse at long intervals)}
\begin{center}\scriptsize
\setlength{\tabcolsep}{3pt}
\renewcommand{\arraystretch}{1.15}
\begin{tabular}{lcccccc}
\toprule
Method & \textbf{CLIP} & \textbf{IR} & \textbf{HPSv2.1} & \textbf{PSNR} & \textbf{SSIM} & \textbf{Speed} \\
\midrule
Cache thường (N=10) & 31.72 & 0.66 & 26.36 & 13.77 & 0.48 & 8.48x \\
Global-Local (N=10) & 31.77 & 0.76 & 27.76 & 14.46 & 0.52 & 8.32x \\
\textbf{\textcolor{wrsBlue}{LearniBridge-100 (N=10)}}\newline{\tiny\color{wrsMuted}(paper method)} & \textbf{\textcolor{wrsBlue}{32.16}} & \textbf{\textcolor{wrsBlue}{0.73}} & \textbf{\textcolor{wrsBlue}{26.68}} & \textbf{\textcolor{wrsBlue}{14.10}} & \textbf{\textcolor{wrsBlue}{0.47}} & \textbf{\textcolor{wrsBlue}{7.96x}} \\
Cache thường (N=25) & 23.08 & -1.48 & 13.71 & 13.12 & 0.42 & 16.87x \\
Global-Local (N=25) & 22.99 & -1.28 & 15.68 & 14.26 & 0.49 & 16.24x \\
\textbf{\textcolor{wrsBlue}{LearniBridge-100 (N=25)}}\newline{\tiny\color{wrsMuted}(paper method)} & \textbf{\textcolor{wrsBlue}{23.09}} & \textbf{\textcolor{wrsBlue}{-1.51}} & \textbf{\textcolor{wrsBlue}{13.72}} & \textbf{\textcolor{wrsBlue}{13.47}} & \textbf{\textcolor{wrsBlue}{0.45}} & \textbf{\textcolor{wrsBlue}{14.77x}} \\
\bottomrule
\end{tabular}
\end{center}
\vspace{0.25em}\wrsTiny{higher is better}
\vspace{0.25em}\wrsQuiet{At 25-step intervals every method collapses (IR below zero, HPSv2.1 near half), so the useful regime in these runs is N=5-10; the qualitative folders show where the remaining failures live.}
\note{Do not read the N=25 rows as a method ranking: all three are far from the original. The qualitative comparison at /Users/nguyenthanhlam/LESA/artifacts/learnibridge\_100\_matched/qualitative\_by\_n is the fair way to inspect the failure modes.}
\end{frame}

\begin{frame}{Qualitative, N=5: Original / Cache / GL / LearniBridge-100}
\wrsFigure[0.99\linewidth]{assets/N5_qualitative.jpg}
\vspace{0.3em}\wrsQuiet{Sheet split in two blocks: rows 1-4 on the left, rows 5-8 on the right. Every row is one prompt, columns are Original, Cache (N=5), GL, LearniBridge-100.}
\note{Source: /Users/nguyenthanhlam/LESA/artifacts/learnibridge\_100\_matched/qualitative\_by\_n/N5.jpg, re-composed as two side-by-side halves so the rows are readable on a slide. Inspect where plain cache drifts and where LearniBridge-100 recovers content.}
\end{frame}

\begin{frame}{Qualitative, N=7: Original / Cache / GL / LearniBridge-100}
\wrsFigure[0.99\linewidth]{assets/N7_qualitative.jpg}
\vspace{0.3em}\wrsQuiet{Sheet split in two blocks: rows 1-4 on the left, rows 5-8 on the right. Every row is one prompt, columns are Original, Cache (N=7), GL, LearniBridge-100.}
\note{Source: /Users/nguyenthanhlam/LESA/artifacts/learnibridge\_100\_matched/qualitative\_by\_n/N7.jpg, re-composed as two side-by-side halves so the rows are readable on a slide. Inspect where plain cache drifts and where LearniBridge-100 recovers content.}
\end{frame}

\begin{frame}{Qualitative, N=10: Original / Cache / GL / LearniBridge-100}
\wrsFigure[0.99\linewidth]{assets/N10_qualitative.jpg}
\vspace{0.3em}\wrsQuiet{Sheet split in two blocks: rows 1-4 on the left, rows 5-8 on the right. Every row is one prompt, columns are Original, Cache (N=10), GL, LearniBridge-100.}
\note{Source: /Users/nguyenthanhlam/LESA/artifacts/learnibridge\_100\_matched/qualitative\_by\_n/N10.jpg, re-composed as two side-by-side halves so the rows are readable on a slide. Inspect where plain cache drifts and where LearniBridge-100 recovers content.}
\end{frame}

\begin{frame}{Qualitative, N=25: Original / Cache / GL / LearniBridge-100}
\wrsFigure[0.99\linewidth]{assets/N25_qualitative.jpg}
\vspace{0.3em}\wrsQuiet{Sheet split in two blocks: rows 1-4 on the left, rows 5-8 on the right. Every row is one prompt, columns are Original, Cache (N=25), GL, LearniBridge-100.}
\note{Source: /Users/nguyenthanhlam/LESA/artifacts/learnibridge\_100\_matched/qualitative\_by\_n/N25.jpg, re-composed as two side-by-side halves so the rows are readable on a slide. Inspect where plain cache drifts and where LearniBridge-100 recovers content.}
\end{frame}

\begin{frame}{Limitations and open questions}
\begin{columns}[T]
\column{0.5\textwidth}
\wrsSection{Current limitations}
\begin{itemize}
  \item [object Object]
  \item Only the final Transformer block is calibrated, so the bridge has to absorb the shift of every bypassed block.
  \item Rank is a design choice. The ablation peaks around r=64 and degrades at larger ranks.
  \item The low-rank and prompt-invariance findings are empirical measurements on the tested models and prompt sets.
\end{itemize}
\column{0.46\textwidth}
\wrsSection{Open questions}
\begin{itemize}
  \item Does the shared low-rank structure survive other backbones, modalities and layer choices?
  \item Can the pre-calibration pass be replaced by data collected during normal inference?
  \item How does the bridge combine with step distillation, quantization or other caching schedules?
\end{itemize}
\end{columns}
\note{End here. No thank-you slide and no leaderboard summary: the mechanism was the point, and the limitations are where the discussion starts.}
\end{frame}
